\section{Models and Training} \label{sec:model}
In our project, we evaluated three different model architectures.
We began with a neural bigram model, which learns to predict the next token based on an embedding of the previous token, using a simple feedforward neural network.
This model serves as a baseline for performance.
We then focused on transformer models and evaluated two variants that differ only in their positional encoding.
The first variant uses sinusoidal positional encodings as introduced in the original Transformer paper \cite{attention}.
The second variant uses a more modern approach, the Rotary Positional Embedding \cite{rope}.
This setup allows us to analyze the effect of the positional encoding choice while keeping the rest of the architecture constant.
Finally, we investigated a gated recurrent unit (GRU) network, which was one of the most widely used sequence modeling architectures prior to the introduction of transformers.
This allows for a comparison between recurrent and attention-based approaches.


\subsection{Bigram} \label{subsec:bigram}
To establish a fundamental baseline for our experiments, we implemented a neural bigram model.
The model predicts the next token based solely on the immediately preceding token, thereby implementing a first-order Markov assumption.
Unlike a traditional count-based statistical model, we constructed this baseline as a simple neural network with learned token embeddings followed by a linear output layer. 
\\
This design choice ensures that the model utilizes the same tokenizer and embedding architecture as our more complex models, enabling a direct and fair comparison of training loss and other evaluation metrics.
By restricting the context window to a single token, this model serves as a meaningful baseline for performance.
It allows us to quantify how much of the text generation task can be solved purely by local token associations, as distinct from the benefits of deeper contextual understanding provided by GRU and Transformer architectures.


\subsection{GRU} \label{subsec:gru}
To overcome the limited context window of the bigram model, we implemented a recurrent neural network (RNN) designed for sequential sequence modeling.
Unlike the bigram model, the GRU processes tokens step-by-step, maintaining a hidden state that captures historical context across the sequence.
\\
We specifically chose the Gated Recurrent Unit (GRU) architecture over the more complex Long Short-Term Memory (LSTM).
GRUs are generally more lightweight, using fewer gates and parameters.
This leads to a simpler implementation, faster training, and easier debugging, while empirical evaluations have shown that they often achieve performance comparable to LSTMs on sequence modeling tasks \cite{chung2014}.


\subsection{Transformer} \label{subsec:transformer}
Our primary focus lies on the Transformer architecture.
Transformers replace recurrence with self-attention mechanisms, enabling each token to directly attend to all previous tokens in the sequence through causal self-attention.
This allows the model to capture long-range dependencies more effectively while also enabling parallelized training.
We chose a decoder-only design for two main reasons.
First, our task is formulated as conditional language modeling: the movie summary and the review are treated as a single token sequence separated by special tokens, making an explicit encoder unnecessary.
Second, this design has established itself as the state-of-the-art for generative text tasks.
While we take the original Transformer from the "Attention Is All You Need" paper \cite{attention} as our starting point, we incorporate modern architectural improvements developed since its publication to enhance performance and stability.
\\
One key improvement we implemented is the use of Pre-Layer Normalization (Pre-LN) instead of the original Post-LN configuration.
We adopted this approach because Pre-LN is known to improve training stability and gradient flow, particularly in decoder-only Transformer architectures.
Although Pre-LN often reduces the need for an explicit learning rate warm-up, we retained a warm-up scheduler in our final setup.
A warm-up phase of 500 steps (corresponding to approximately $1/16$ of an epoch, depending on model size) was used, as warm-up is commonly employed to stabilize optimization during the early stages of training. 
\footnote{The warm-up scheduler was initially introduced during a debugging phase when the model failed to reduce training loss.
The root cause was later identified as a data preprocessing error in which input tokens were inadvertently shuffled.
This issue was diagnosed using a single-batch overfitting test.
Although the underlying bug was resolved, the warm-up mechanism was retained in the final pipeline as a conservative measure to guard against potential initialization instabilities.}
To optimize results, we also evaluated two different positional encoding schemes: Sinusoidal Positional Embeddings and Rotary Positional Embeddings (RoPE), allowing us to compare additive versus rotary-based position encodings.


\subsection{Tokenizer} \label{subsec:tokenizer}
To get a functional MVP early in the project, we initially employed a simple character-level tokenizer in combination with the Bigram model.
The character-level tokenizer was used exclusively for this early Bigram-based prototype and is not part of any subsequent experiments or reported results.
\\
For all final models and experiments, we employed a byte-level Byte Pair Encoding (BPE) tokenizer with a fixed vocabulary size of 4{,}000 tokens.
This vocabulary size was chosen as a trade-off between model efficiency and the linguistic expressiveness required to capture the diverse vocabulary of film reviews.
In addition to standard control tokens such as [BOS] (Beginning of Sequence), [EOS] (End of Sequence), and [PAD] (Padding), we introduced task-specific delimiters to structure the input for conditional generation.
Specifically, we defined a [SYN] token to mark the start of the movie summary and a [REV] token to signal the beginning of the generated review.
\\
We chose a byte-level BPE tokenizer for several reasons.
Operating at the byte level ensures complete coverage of all UTF-8 encoded characters, including emojis, special symbols, and informal language commonly found in film reviews, thereby eliminating the possibility of out-of-vocabulary tokens.
Furthermore, BPE provides a flexible compromise between character-level and word-level representations, allowing the tokenizer to learn subword units based on the frequency distributions observed in the training corpus, thus adapting to the characteristic vocabulary and style of film reviews.
Although we initially implemented the BPE algorithm from scratch to gain a deeper understanding of the algorithm, the final training pipeline utilizes the Hugging Face Tokenizers library.
This decision was driven by computational efficiency, as the Rust-backed implementation offers significantly faster tokenization speeds during data preprocessing compared to a native Python implementation.


\subsection{Training} \label{subsec:training}
The entire training pipeline was implemented in PyTorch. All models were trained using a standard next-token prediction objective.
To ensure comparability across architectures, we applied consistent training procedures and evaluation protocols across all models, with adjustments made only where necessitated by architectural constraints.

\paragraph{Data Splitting and Loading}
We split the dataset into 80\% training, 10\% validation, and 10\% test data.
To maximize computational efficiency, all reviews were concatenated into a continuous token stream separated by [BOS] and [EOS] tokens, from which we sampled random fixed-length chunks during training.
This streaming approach eliminates padding overhead and naturally packs multiple short reviews into a single batch.

\paragraph{Optimization}
For the GRU and Transformer models, we utilized the AdamW optimizer with a weight decay of $0.01$ to mitigate overfitting and a learning rate of $3 \times 10^{-4}$.
In contrast, the Bigram was trained without weight decay and with a higher learning rate of $1 \times 10^{-3}$.
We made these adjustments to account for the Bigram's significantly smaller parameter space, which requires less regularization and converges more rapidly.
\\
During training, model performance was monitored at regular intervals.
In the first 1000 training steps, evaluations were performed every 100 steps to closely track early training dynamics.
After this initial phase, evaluations were conducted every 500 steps.
At each evaluation checkpoint, we computed the average loss over multiple independently sampled batches from both the training and validation sets. 
Specifically, we evaluated 50 batches for the Transformer and GRU models and 30 batches for the Bigram baseline.
Averaging over multiple batches reduces the variance in single-batch estimates and yields more stable and reliable performance measurements.
During evaluation, models were temporarily switched to evaluation mode to disable stochastic components such as dropout.
The validation loss served as the primary criterion for early stopping.
Training was terminated if the validation loss failed to improve by at least a minimum delta of $10^{-3}$ over five consecutive evaluation checkpoints.
This criterion indicates convergence or the onset of overfitting and ensures that all models are trained until a comparable level of performance stability is reached.
\\
While initial experiments were conducted on a local MacBook with an Apple M1 processor, final training runs were executed on GPU-equipped systems (AMD RX 6650XT) to significantly reduce training time for the GRU and Transformer architectures.
